% !TeX root = ../main.tex
\chapter{Reinforcement-Learning Teacher and Automated Demonstrations}
\label{ch:teacher}
\section{Why Reinforcement Learning Is Used}
The teacher is needed because the student requires action labels and the
proposed workflow seeks to obtain them without manual driving demonstrations.
RL supplies a mechanism for learning a policy from simulated interaction:
steering changes the trajectory, and reward feedback evaluates the subsequent
motion. Forward travel, clearance and collision outcomes can be computed by
the virtual environment without an operator specifying an action for every
scan. The resulting teacher can then generate the supervised dataset.

The distinction between reward design and action labelling is central. A reward
states what outcomes are desirable; it does not enumerate the steering command
for every encountered state. RL is useful here because a locally plausible turn
can affect later collision risk, so the training objective must consider a
sequence of decisions. This does not remove the need to design suitable rewards
or to qualify the teacher before collecting its demonstrations.

A manually programmed expert or a classical planner could also generate labels.
Such a teacher would require an explicit control strategy and appropriate
vehicle/task assumptions. The present project selects RL to learn that strategy
from interaction. It has not compared the cost or quality of all possible
teacher sources, and it does not claim that RL is the only route to automated
supervision.

\section{Why PPO Is Selected}
PPO is an on-policy actor--critic algorithm that updates a policy using recently
collected rollouts. Its clipped surrogate limits the incentive for some large
probability-ratio changes, providing a practical update mechanism without
requiring gradients through the simulator~\cite{schulman2017ppo}. The chosen
implementation provides the policy, value estimation and rollout machinery
within a consistent training interface~\cite{raffin2021sb3}.

These properties suit this workflow: the simulator can repeatedly reset and
execute the fifteen-action interface, while the host performs optimisation.
The selected run uses four parallel environments and 512 rollout steps per
environment before an update. PPO is an engineering choice supported by the
available implementation and observed teacher qualification, not a demonstrated
winner over every RL algorithm. On-policy learning also requires new interaction
after updates, so its simulation cost remains part of the training cost.
No algorithm-comparison experiment establishes that PPO minimises that cost.

The alternatives clarify the choice. Deep Q-learning learns action values and
has been demonstrated with discrete actions~\cite{mnih2013playing}; the
fifteen-action interface therefore does not uniquely require PPO. TRPO uses a
trust-region approach to policy optimisation~\cite{schulman2015trpo}, while
PPO offers a clipped first-order surrogate that is simpler to incorporate into
the selected rollout workflow~\cite{schulman2017ppo}. The original Soft
Actor-Critic formulation combines off-policy learning and entropy maximisation
for continuous-control benchmarks~\cite{haarnoja2018sac}; using that formulation
would require reconsidering the action interface. Discrete SAC variants are
not ruled out by this observation. The project prioritises a reproducible
categorical policy and an established implementation, while accepting PPO's
interaction cost. These are design reasons, not measured advantages over DQN,
TRPO or SAC on the present maps.

\section{Which Model the Teacher Uses}
The teacher is Stable-Baselines3's \texttt{MlpPolicy}, an
\texttt{ActorCriticPolicy} with separate feed-forward policy and value branches.
Inspection of the frozen checkpoint confirms the dimensions in
Table~\ref{tab:teacher_network}. Both branches use two 128-unit hidden layers
with Tanh activations. The policy head outputs fifteen logits defining a
categorical action distribution; the value head outputs one scalar estimate
of discounted return. The critic supports training and does not provide the
student's action labels.

\begin{table}[htbp]
\centering\small
\caption{Neural architecture of the frozen teacher, verified from the saved model.}
\label{tab:teacher_network}
\begin{tabularx}{\textwidth}{l l >{\raggedright\arraybackslash}X}
\toprule Branch & Layer widths & Output role \\\midrule
Actor & $292\to128\to128\to15$ & Categorical distribution over steering commands. \\
Critic & $292\to128\to128\to1$ & Scalar value estimate used during learning. \\
\bottomrule
\end{tabularx}
\end{table}

An MLP is a straightforward function approximator for the fixed-length numeric
observation used here. The implementation flattens the distance field and adds
three state features; it does not process camera images or use a recurrent
history. A convolutional or recurrent policy could exploit different structure,
but no comparison establishes a benefit for this task. The MLP is therefore a
concrete baseline, not a claim that spatial structure or memory are unimportant.

Counting weights and biases gives
$(292+1)128+(128+1)128+(128+1)15=55{,}951$ actor parameters and
54,145 critic parameters. FP32 parameter storage is 223,804 bytes (218.6~KiB)
for the actor and 440,384 bytes (430.1~KiB) for both branches, where
$1~\mathrm{KiB}=1024$ bytes. These derived sizes exclude activations, runtime
code and training state. Actor-only inference needs neither critic nor optimiser;
parameter storage alone does not establish board feasibility.

The teacher uses privileged geometry because its immediate purpose is to learn
useful driving behaviour and provide supervision. This choice separates
learning the task from reconstructing that geometry from an imperfect scan,
in the spirit of privileged-to-sensor transfer~\cite{chen2019cheating}. It also
creates an information gap that imitation may not fully overcome. Directly
installing the teacher on the vehicle would not solve that gap: its map-derived
input is unavailable in the proposed LiDAR-only onboard interface, regardless
of whether the actor's weights could fit in memory.

\section{Teacher Observation and Learning Objective}
The teacher observes a local map-derived distance field on a $17\times17$
grid, together with forward speed, yaw rate and previous action. Its input
therefore has $289+3=292$ components. The grid spans 0.4~m behind to 1.4~m
in front of the vehicle and extends 0.9~m to each side; distance values are
clipped at 0.5~m according to the stored observation configuration. It is a
privileged representation because it is computed from simulator geometry and
state rather than reconstructed from the student's scan.

Let $o_t^T$ denote this observation and $\pi_\theta(a_t\mid o_t^T)$ the teacher
policy. Training uses Stable-Baselines3~\cite{raffin2021sb3} with the surrogate
\begin{equation}
 L^{\mathrm{clip}}(\theta)=\mathbb{E}_t\!\left[
 \min\!\left(r_t(\theta)\hat A_t,
 \operatorname{clip}(r_t(\theta),1-\epsilon,1+\epsilon)\hat A_t\right)\right],
\end{equation}
where $r_t$ is the new-to-old action probability ratio and $\hat A_t$ is an
advantage estimate~\cite{schulman2017ppo}. Generalized advantage estimation
(GAE) provides a value-based estimator with a bias--variance trade-off~\cite{schulman2015gae}.
The stored configuration uses
$\epsilon=0.2$, learning rate $3\times10^{-4}$, discount factor 0.99,
GAE parameter 0.95, batch size 256 and ten update epochs. Two hidden layers
of 128 units are specified, with entropy coefficient 0.005 and value-loss
coefficient 0.5. These settings describe the selected run, not a new
hyperparameter search conducted for this draft. The saved teacher artifacts
are identified in Appendix~\ref{sec:evidence_sources}.

\section{Reward and Termination}
The reward combines bounded forward travel, speed tracking, geometric safety,
steering smoothness and terminal failure penalties. Table~\ref{tab:reward}
summarises the coefficients in the saved teacher configuration. For example,
with geometric clearance $c_t$, the near-obstacle term is
\begin{equation}
 r_t^{\mathrm{near}}=-0.5\max(0,1-c_t/0.25)\,\Delta t.
\end{equation}
The forward-risk term uses the square of a clipped deficit relative to
0.45~m. In the privileged mode these safety terms use map geometry, not
student LiDAR values. Distances in the expressions are in metres.

\begin{table}[htbp]
\centering\small
\caption[Teacher reward components]{Recorded teacher reward components. Time-scaled terms use $\Delta t$.}
\label{tab:reward}
\begin{tabularx}{\textwidth}{p{3.2cm}X}
\toprule Component & Implementation summary \\\midrule
Travel & $+10$ per metre of bounded forward travel; suppressed on failure. \\
Speed error & $-3|v-v_0|\Delta t$. \\
Near obstacle & Linear deficit below 0.25~m, coefficient 0.5 per second. \\
Forward risk & Squared deficit below 0.45~m, coefficient 2 per second. \\
Steering change & $-0.1|a_t-a_{t-1}|/14$. \\
Turning & $-0.15((a_t-7)/7)^2\Delta t$. \\
Failure & $-20$ for the applicable collision or motion-failure category. \\
\bottomrule
\end{tabularx}
\end{table}

Collision, insufficient travel, small repeated loops, persistent low speed
and leaving the workspace are failure conditions. The stored stuck monitor
uses a 3~s window and a 0.02~m travel threshold. Loop monitoring uses a 12~s
window with additional duration, angle and spatial-span conditions. These
checks distinguish continued forward cruising from remaining near one place.
Reaching the configured time limit without a failure is treated as episode
completion, not arrival at a navigation goal.

\section{Training Provenance and Teacher Qualification}
The collection teacher is the checkpoint at 2,392,588 training steps in the
second training stage. That stage continued from a 42,588-step parent and
ended at 2,500,520 steps. Later multi-map and dead-end runs exist in the
project but are not the source of the reported demonstrations. The presence
of a multi-map training script therefore does not establish that the selected
teacher was trained on multiple maps.

Figure~\ref{fig:teachercurve} shows the archived second-stage learning curve.
Its selected checkpoint is distinguished from the end of the run. The frozen
teacher's separate qualification summary records eight episodes on T3,
approximately 4,800 simulation seconds and 959.90~m of travel, with no recorded
collision or listed motion failure. This establishes behaviour under that
qualification protocol. Successful collection on other maps is additional
operational evidence, but is not presented as an independent four-map teacher
qualification experiment.

\begin{figure}[htbp]
\centering\includegraphics[width=\textwidth]{teacher_training_curve.pdf}
\caption[Teacher training curve and selected checkpoint]{Archived second-stage teacher training curve and selected checkpoint.
The evaluation-reward axis has a narrow range; small fluctuations should not
be interpreted as large changes in navigation capability.}
\label{fig:teachercurve}
\end{figure}

\section{Observation--Action Pairing}
After qualification, the teacher is frozen. The collector uses deterministic
prediction, selecting
\begin{equation}
 a_t^T=\arg\max_a\pi_\theta(a\mid o_t^T),\qquad
 \mathcal D=\{(x_t,a_t^T)\}_{t=1}^{N}.
\end{equation}
No person supplies the per-sample action labels in this stage. The forest
later fits $\mathcal D$ by supervised classification; it receives neither
PPO updates nor simulator rewards. This is behaviour cloning from hard action
labels, rather than direct RL training of tree parameters.

A demonstration sample is $(x_t,a_t^T)$, where $x_t$ is the valid LiDAR vector
at the decision boundary and $a_t^T$ is the teacher command subsequently
executed. The collector checks scan freshness instead of assuming that all
sensor and physics callbacks are simultaneous. Missing decision scans are
skipped; previous scans are not duplicated to manufacture a row. Five
consecutive misses cause collection to abort under the recorded configuration.

Each published batch includes teacher and map fingerprints, configuration,
seeds, episode outcomes, and inclusive row ranges linking episodes to files.
Training files contain 200 distance values and one action label. Timing logs
and runtime metadata are retained separately from the CSV copies used for
student learning. This division supports both compact training input and
an audit of the collection process.

\section{Starting Conditions and Dropout Augmentation}
Ordinary teacher starts use position perturbations up to 0.03~m and yaw
perturbations up to $10^\circ$. Near-wall starts use additional geometric
constraints, including 0.04--0.08~m body clearance, $10$--$20^\circ$ yaw offsets
and 0.15--0.40~m forward clearance. Candidate starts are checked before use.
These perturbations expose the teacher to corrective situations, but do not
establish student recovery from the same situations.

For each original scan, offline dropout chooses approximately 5\% of its
nonzero measurements and sets them to zero without changing the action label.
The count is rounded as $\lfloor0.05n_{\mathrm{nz}}+0.5\rfloor$. The original
and augmented rows retain a common map assignment. Their labels are identical,
so augmentation increases the number of training rows without doubling the
number of independently visited states.

\section{Dataset Composition and Later Batch}
Table~\ref{tab:batches} summarises the two batches used in the current analysis.
The batch metadata are identified in Appendix~\ref{sec:evidence_sources}.
The October 23 batch supplies student training and validation. The October 28
batch is used only to evaluate the frozen models on new teacher trajectories
on the same maps. Neither batch contains student-driven recovery demonstrations.

\begin{table}[htbp]
\centering\small
\caption[Demonstration batch composition]{Demonstration batches; each episode is limited to 3,600 simulation seconds.}
\label{tab:batches}
\begin{tabular}{lrr}
\toprule Quantity & 23 October & 28 October \\\midrule
Maps & 4 & 4 \\
Episodes per map & 20 & 15 \\
Completed episodes & 80/80 & 60/60 \\
Original rows & 1,439,983 & 1,079,979 \\
Dropout rows & 1,439,983 & 1,079,979 \\
Skipped decision scans & 17 & 21 \\
Ordinary / near-wall starts & 60 / 20 & 44 / 16 \\
\bottomrule
\end{tabular}
\end{table}

The original demonstration batch contains only labels 5, 6, 7, 8, 9 and 12. Label 7 accounts
for 873,837 original rows, approximately 60.68\%, whereas label 12 occurs only
22 times. The later batch also contains only these six labels, with label 12
occurring 12 times. Consequently, the fifteen-command interface should not be
confused with a dataset that supplies evidence for all fifteen actions.
Long runs provide many temporally adjacent samples, not an equal number of
independent driving situations. This limitation motivates per-map and
per-episode reporting in the subsequent experiments.
